\documentclass[10pt,twocolumn]{article}
\usepackage[letterpaper,margin=0.76in,columnsep=0.26in]{geometry}
\usepackage[T1]{fontenc}
\usepackage{lmodern,microtype,amsmath,amssymb,booktabs,graphicx}
\usepackage{xcolor}
\usepackage{xurl}
\usepackage[colorlinks=true,allcolors=blue!45!black]{hyperref}
\hypersetup{pdftitle={Measuring the Value of Failed Code in Budgeted Model Recovery},pdfsubject={Audited recovery study; incomplete descriptive MBPP+ bank}}
\usepackage{enumitem}
\setlist{nosep,leftmargin=*}
\setlength{\parskip}{2pt}
\renewcommand{\topfraction}{0.92}
\renewcommand{\dbltopfraction}{0.92}
\renewcommand{\textfraction}{0.06}
\renewcommand{\floatpagefraction}{0.75}
\renewcommand{\dblfloatpagefraction}{0.75}
\newcommand{\Usd}{\mathrm{USD}}
\input{generated/results-macros.tex}
\title{Measuring the Value of Failed Code in Budgeted Model Recovery}
\author{}
\date{Research draft, 22 October 2026}
\begin{document}
\maketitle
\begin{abstract}
After a coding attempt fails, a downstream model can solve the task again,
reuse the failed program, inspect public test feedback, or combine these inputs.
The value of the program to the solver need not equal its value to a controller
choosing the next action. We study both quantities with matched, independently
sampled recovery requests and frozen, source-trained selectors. An initial
HumanEval+ producer-transfer study does not establish an advantage over a strong
development-selected static recovery rule. We therefore freeze an external
MBPP+ replication that adds a feedback-only prediction control, direct
model-by-context interactions, conservative billing, and an outcome-blind
semantic forecast baseline. Among 94 complete external failure states, adding
failed code to feedback reduces 9B success by 22.0 percentage points and 27B
success by 3.9 points. The direct model interaction is $-18.1$ points, with a
nominal family-adjusted 99.583\% interval $[-31.6,-5.3]$. However, the frozen
artifact-feature selector's primary utility advantage over feedback-only is
\$0.000111 per failure, with a descriptive 95\% interval spanning zero.
Two permanently missing solver responses prevent the planned complete-bank
confirmatory analysis. These results support a reproducible measurement study;
they do not establish a new routing method or a validated agent controller.
\end{abstract}

\section{Introduction}

A failed program is an ambiguous input to the next coding attempt. It can
contain a useful interface, a nearly correct algorithm, or an attractive
implementation of the wrong interpretation. Test feedback can clarify the
mistake without requiring that implementation to be preserved. A controller
that observes both inputs must consequently answer two different questions:
which information should the next solver receive, and which information helps
predict that choice?

These questions matter when choosing whether to pay for another attempt.
Code reuse adds input tokens and can constrain the next response. A forecaster
adds a separate inference charge. A controller therefore needs to improve the
verified outcome enough to offset both the chosen action's cost and its own
overhead. Better average probability calibration does not by itself establish
better decisions: action selection depends on differences between probabilities
and costs.

Recovery routing is already studied by CodeRescue~\cite{coderescue}, and
packet decomposition with placebo controls is central to Falsification, Not
Exposure~\cite{falsification}. Our aim is narrower than proposing another
reflect/replan/escalate interface. We measure the added value of artifact
information relative to verifier feedback, distinguish context selection from
model selection, and ask whether frozen predictors transfer to a different
benchmark. We publish all matched observations, forecasts, billing reservations
and the replay needed to inspect those claims.

Our initial HumanEval+ evaluation exposes a baseline-selection issue. A state
selector improves utility against a static action selected using single-draw
priors, but its advantage over a static action selected using the actual
two-draw replay objective is uncertain. This motivates a stricter external
test with a feedback-only predictor and a static comparator chosen for the
same finite decision horizon. The earlier held-out failures become explicit
training data for that test; they are not counted again as fresh evaluation.

\section{Related work}

CodeRescue trains a post-failure recovery router and calibrates its budget
penalty~\cite{coderescue}. Its disclosed routing idea remains prior art despite
the preprint's administrative withdrawal; our experiments do not reproduce its
trained system. Falsification, Not Exposure decomposes repair inputs and
compares them with resampling and shape-matched placebos~\cite{falsification}.
Our code/feedback factorial and irrelevant-context control are therefore not
novel contributions. DARC freezes diagnosis-guided admissible recovery
interventions from development failures~\cite{darc}. DDBench compares bounded
debugging context with symptom-only requests in distributed-system
repair~\cite{ddbench}. These settings differ from our short-function
checkpoint bank, but already show why generic context expansion is not a new
research premise.

Self-repair and self-debugging already examine the use of failed code and
execution feedback~\cite{selfrepair,selfdebug}. Re-FORC predicts future reasoning
reward and cost for decisions that include stopping and model
selection~\cite{reforc}; AI Functions documents economic routing and forecasting
interfaces~\cite{aifunctions}. We compare limited local and API forecasting
baselines rather than claiming these general decision ideas as new.
HumanEval+ and MBPP+ provide the executable checks in our
study~\cite{evalplus,mbpp}. They are small-program benchmarks, not substitutes
for repository-level agent evaluation.

\section{Problem and estimands}

Let a visible failure state be $s=(x,a,f)$: task text $x$, failed artifact $a$,
and public verifier feedback $f$. A next-solver action is a model $m$ and context
condition $c$. The selectable conditions are fresh task text, artifact only,
feedback only, and both. An irrelevant artifact-plus-feedback package is a
diagnostic control.

For a fixed state, let $p_{m,c}(s)$ be the probability that a requested action
returns a complete solution passing both the public and evaluation-only test suites. One relevant
quantity is the incremental artifact effect given feedback,
\[
  \delta_{a\mid f,m}(s)
    =p_{m,\mathrm{both}}(s)-p_{m,\mathrm{feedback}}(s).
\]
The comparison with fresh generation answers a different operational question:
it includes the value of public feedback as well as that of the artifact.
Neither comparison identifies a cognitive mechanism such as anchoring.

The selector observes a specified view of $s$ and estimates success for each
action. With declared solution value $v$, it chooses an admissible action to
maximize $v\widehat p_{m,c}(s)-\widehat k_{m,c}$. Costs used for this choice come
from source training observations. Actual evaluation receipts determine the
reported cost, while worst-case token reservations determine advance admission.
The policy abstains when no admissible action has positive predicted net value.

Our primary contrast asks whether adding artifact features to an otherwise
identical feedback view improves the utility of the resulting frozen selector.
This differs from asking whether supplying the artifact to the \emph{solver}
improves generation. A controller may profit from inspecting an artifact and
then discard it, or reuse an artifact without deriving any benefit from artifact
features in its own decision model.

\section{Experimental design}

\subsection{Benchmarks and visible information}

We pin EvalPlus 0.3.1, HumanEval+ v0.1.10 and MBPP+ v0.2.0 by dataset checksums.
The public channel explicitly exposes the official base test inputs and
acceptable reference outputs, at most three failing examples, the failed
fraction and verdict. Plus tests remain evaluation-only and are absent from the
public-verifier subprocess payload and all model requests. Full success requires
base and Plus success. A replay stops on a public pass, including a public false
positive later rejected by Plus.

Canonical references undergo the same public and full-suite verification before
model collection. HumanEval/47 has a contradictory public example and
HumanEval/32 is incompatible with the pinned checker's special success path;
both are excluded from the source study. For MBPP+, all 378 base reference
checks pass and 377 full-suite checks pass. Mbpp/599's official sum-of-range
implementation exceeds the isolated full-suite verifier limit, so it is excluded
before model responses. A separate public-prompt assertion audit finds all
378 examples compatible with their references when the test harness supplies
\texttt{math}. That check does not prove every English specification correct.

Generated code runs in macOS OS sandboxes without network, credentials,
private user files or access to the evidence bank. Public and full-suite payloads
remain separate. MBPP input transport uses EvalPlus's task-specific serializers,
including complex values, and its official special-oracle rules.

\subsection{Source study and transfer cohort}

The earlier source study generated 83 development prefixes with Gemma 3 4B IT
and 60 evaluation prefixes with Ministral 3B 2512, after a disclosed producer
availability amendment. It observed 26 development and 25 evaluation failures.
The evaluation was consequently a Gemma-to-Ministral producer-transfer test,
not a completed same-producer validation. These 51 failed states are now all
source training data for the new MBPP+ experiment.

For MBPP+, we numerically sort the 377 eligible task IDs, shuffle them once with
seed 20261025 and generate one Ministral prefix per task. Prefix parameters are
temperature 0.6 and a 1536-token output ceiling, with a user-only JSON solve
instruction. The producer is pinned to Mistral's endpoint, whose precision is
reported as unknown. Only a public failure triggers recovery; format, provider,
syntax and ordinary wrong-answer failures are retained.

The first at most 96 failures in the frozen shuffled order form the recovery
cohort. We do not select by Plus scores, replace unhelpful artifacts or collect
until a significance threshold is crossed. One late prefix transport failure
after 347 observed prefixes has no preserved response identity. All 96 selected
failures precede it, so its unobserved outcome cannot change the cohort.
The continuation keeps its full reservation, does not reissue that task or any
observed request, and completes only unadmitted prefixes. It yields 376 observed
prefixes and one censored prefix among 377 eligible tasks.

\subsection{Matched recovery requests}

Each selected failure receives three independent draws in five context
conditions for Qwen3.5 9B and 27B. Both consumers are pinned to
\texttt{siliconflow/fp8}, with reasoning and fallbacks disabled and no provider
sampler seed. They use the same solve instruction, temperature 0.6 and 1536-token
output ceiling. This is a practical comparison of two nominal model capacities,
not an identified causal effect of parameter count or intelligence.

The five conditions are fresh, code, feedback, both, and distractor. The
distractor is a different HumanEval+ training failure, chosen deterministically.
Its code and feedback fields are independently truncated or padded to the
relevant fields' character lengths. This is \emph{not token matching}; we report
input-token receipts and retain that limitation when interpreting relevance.
Distractors are never selectable actions.

Submission order is randomized within each checkpoint. Four bounded network
workers collect responses; two isolated verifier workers score recovery
candidates. Prefixes and canonical reference audits are serial. We admit a
complete 30-request checkpoint block only if its sum of worst-case reservations
fits the uncommitted collection budget. No automatic inference retry or provider
substitution is allowed after uncertain transport.

Two subsequent recovery requests have no preserved response identity: a 27B
code-only third draw on Mbpp/426 and a 9B fresh first draw on Mbpp/84. Neither
request is repeated and both full reservations remain held. After the first
recovery stop and a disclosed descriptive audit of 47 complete states, the user
explicitly authorizes only the 1,466 planned requests never admitted before that
stop. A later 433-request segment is a strict subset of this same allowlist,
with source hashes, earlier lineages and journal prefixes preserved. Models,
prompts, weights, verifier and estimands do not change. This exception does not
restore the original complete-bank design: at most 94/96 selected states can
have all 30 draws. All complete-case intervals are consequently descriptive.
We also retain all 96 selected states in worst-case missing-state utility bounds;
these bounds are not confidence intervals and do not impute missing responses.

\subsection{Frozen local predictors and policies}

Source fitting uses Beta(1,1) action priors, ridge penalty 10, source-only feature
normalization and clipping to $[0.01,0.99]$, without target fitting or a feature
search. The task view contains log prompt length and line count. The feedback
view adds the public failure fraction and NameError, SyntaxError, TypeError,
timeout and format indicators. The state view adds exactly three artifact
features: log code length, log AST size and syntax validity. A global prior is an
additional control. These structural features do not exhaust semantic artifact
information.

Weights and training-source hashes are committed before any MBPP response.
All selected-state predictions are saved before the first target recovery
request. The verified-solution value is fixed at \$0.01 and task ceilings are
\$0.003, \$0.01 and \$0.02. A policy uses at most two independent draws of its
chosen fixed action from the original checkpoint; it does not update from the
first child's code or feedback.

The strong static comparator is selected by complete two-draw replay on the
51 source failures separately for each budget. Its choices are small-fresh at
\$0.003 and large-fresh at \$0.01 and \$0.02. A model-only ablation selects the
consumer while always supplying both code and feedback. All eight fixed actions
are reported. Conditional fresh-model controls without a producer charge are
diagnostics for this failed-prefix cohort, not whole-benchmark deployment
strategies; passing-prefix direct controls are not collected in the replica.

An additional constant escalation diagnostic, added after recovery collection
began and before treatment statistics were computed, tries 9B fresh and then
27B fresh after public failure. Both calls start from the original task, without
a new child artifact handoff. It consumes recorded draws, not additional model
calls, and does not replace the primary or its frozen comparisons.

\subsection{Semantic forecast control}

Before any recovery response, we additionally fix a Qwen 9B forecaster with
task-only, feedback-only and full-state views. All views receive identical
source action priors, source mean costs, model identities and forecast
instructions. Only the full-state view receives the failed program. No Plus
datum or recovery outcome enters its request.

To avoid overlapping budget allowances, these API forecasts are collected only
after the primary bank is finished, with their prompt and code already fixed.
They are a prespecified outcome-blind control, \emph{not physically prospective
forecast responses}. Invalid outputs retain their fee and use the declared
0.5 probability fallback. In replay the forecast must first be admitted under
its worst-case reservation and its measured fee is charged before solver
selection.

\subsection{Statistics and accounting}

The primary statistic is state-minus-feedback policy mean net utility at the
\$0.01 task ceiling, using draws 0 and 1. We bootstrap independent task clusters
10,000 times, seed 20261025, and report a two-sided 95\% percentile interval.
Success and cost differences accompany utility. An interval spanning zero is not
evidence of equivalence. Resampling is over target tasks only, conditional on
the saved source-trained weights; it does not propagate source-training
uncertainty or refit a predictor on target outcomes.

Prespecified secondary context contrasts are code-fresh, both-fresh,
both-feedback and both-distractor for each model, plus the four corresponding
small-minus-large interactions. All twelve use Bonferroni-adjusted nominal
99.583\% intervals as well as ordinary 95\% intervals. The interaction family was
fixed before any target response; separate significance results for two models
would not themselves establish an interaction.

We also report Brier scores, squared errors for predicted both-feedback and
both-fresh differences, and replay sensitivity to draw orders (1,2) and (2,0).
Observed three-draw effects remain noisy. An exploratory cross-draw-product
variance diagnostic estimates between-task benefit variation without treating
all sampling noise as signal; negative finite-sample estimates remain
unclipped. It assumes conditionally independent provider draws and is not a new
estimator or proof of predictability.

Every inference admission reserves the conservative input/output token ceiling
before network activity. Known charges use provider receipts. Unknown billing
retains the full reservation. Explicit, identified empty provider errors use a
disclosed policy-backed zero only when covered by OpenRouter's published
insurance~\cite{openrouter}; that is not a token receipt. Training-data collection,
evaluation collection and forecaster fees are disclosed separately. Local
verifier elapsed time is not converted into a fabricated dollar charge.
The approved additional campaign ceiling is \$2. Prior charges and reservations
contribute \$0.5149634; the separate primary replication cap is \$1.48, giving a
combined upper allowance of \$1.9949634 before optional semantic calls. The latter
may use at most \$0.08 and must additionally fit the remaining campaign balance
after all known charges and outstanding reservations. No continuation receives
a fresh allowance.

\section{Results}

\subsection{Earlier HumanEval+ transfer result}

\begin{table}[t]
\centering
\caption{Single-draw full-suite success on 25 previously held-out HumanEval+
failures. Two draws per cell; all 51 source failures subsequently become training
data for MBPP+.}
\label{tab:source}
\small
\begin{tabular}{lrrrrr}
\toprule
Model & Fresh & Code & Feedback & Both & Irrel.\\
\midrule
9B & 66\% & 50\% & 64\% & 62\% & 64\%\\
27B & 84\% & 86\% & 80\% & 86\% & 86\%\\
\bottomrule
\end{tabular}
\end{table}

On the earlier 25 failed evaluation tasks, code-only recovery changes 9B success
by $-16$ percentage points relative to fresh generation, with an exploratory,
unadjusted 95\% interval $[-30,-4]$. The corresponding 27B point difference is
$+2$ points. These separate intervals do not establish a cross-model interaction,
nor do they justify a universal rule to discard failed code.

Across all 60 original evaluation tasks, the state selector and strong
source-selected static rule each solve 55/60 tasks at the \$0.01 ceiling.
Mean API costs are \$0.00037164 and \$0.00040016, respectively. Their utility
difference is \$0.00002851 per task, with 95\% interval
$[-\$0.00046099,+\$0.00054027]$. Against the original static action chosen using
single-draw priors, the utility difference is positive
(\$0.00040765, interval $[\$0.00003482,\$0.00095225]$).
Both comparisons are retained: the apparent improvement depends materially on
matching the static comparator to the decision horizon.

Benefit prediction is not established in that study. State-minus-task Brier
and observed-benefit-MSE differences both have intervals spanning zero. A
model-only ablation solves 56/60 tasks, so the joint selector does not establish
added value from context choice over model routing. These findings motivate the
new feedback-only control and external evaluation; they are not refit into a
positive source evaluation claim.

\subsection{Frozen MBPP+ replication}
\input{generated/bank-status.tex}

\begin{table}[t]
\centering\small
\caption{Conditional single-draw full-suite success on MBPP+ failed tasks.
Three independent draws per model/context cell. Rates use the same complete
cases across all ten cells; missing observations are retained in the archive.}
\label{tab:mbpp-context}
\input{generated/context-rates.tex}
\end{table}

\begin{table}[t]
\centering\small
\caption{Two-draw MBPP+ replay at the \$0.01 task ceiling. Cost and utility are
USD $\times1{,}000$ per failed task, including the producer. Escalation* is the
additional diagnostic added during collection.}
\label{tab:mbpp-policy}
\input{generated/policy-results.tex}
\end{table}

\begin{figure*}[!t]
\centering
\includegraphics[width=\textwidth]{evidence-audit/replication.pdf}
\caption{Descriptive MBPP+ complete-case comparisons. Top: context effects with
nominal 99.583\% task-bootstrap intervals for the twelve-comparison family;
the four direct model interactions are in Table~\ref{tab:all-context-effects}.
Bottom: two-draw recovery success and mean API cost on the same failed tasks,
with producer cost charged. The scatter does not replace a paired utility
comparison. Both panels exclude the same incomplete states; all-selected-state
bounds remain explicit in the text.}
\label{fig:mbpp}
\end{figure*}

\begin{figure}[!t]
\centering
\includegraphics[width=\columnwidth]{evidence-audit/policy-utility.pdf}
\caption{Paired state-minus-control mean utility on MBPP+ at the \$0.01
ceiling. The feedback comparison is primary; others are secondary. Intervals
are descriptive 95\% task-bootstrap intervals. Escalation* was added during
collection and starts both models from the original task.}
\label{fig:paired-utility}
\end{figure}

\begin{table}[t]
\centering\small
\caption{Frozen local predictor scores on MBPP+ complete cases.
Benefit MSE averages both-minus-feedback and both-minus-fresh losses for
each model and includes three-draw outcome noise.}
\label{tab:mbpp-prediction}
\input{generated/prediction-results.tex}
\end{table}

All authorized primary requests are accounted for: 2,878 observed recovery
responses and two censored cells among the 2,880 planned requests. Complete-case
comparisons use 94 states, three draws per model/context cell, and therefore
282 observations per reported success rate (Table~\ref{tab:mbpp-context}). The
40 public false positives among initial responses and 360 among recovery
responses remain failures under the full-suite criterion. Three identified
empty provider errors are retained rather than dropped as missing generations.

For 9B, code alone reduces success from 48.6\% fresh to 26.6\%; both code and
feedback reach 30.9\% versus 52.8\% with feedback alone. For 27B, the analogous
rates are 61.3\%, 50.4\%, 58.5\% and 62.4\%. The both-minus-feedback effect is
$-22.0$ points for 9B (nominal adjusted interval $[-33.7,-11.0]$) and $-3.9$
points for 27B ($[-14.2,+5.3]$). Their direct difference is $-18.1$ points
($[-31.6,-5.3]$); this interaction, rather than separate significance labels,
supports different context effects in this particular comparison.
Code-only versus fresh is negative for both models, including the larger one.
The results consequently do not imply that larger models always benefit from
failed code. Both-versus-irrelevant is also negative for 9B, but the
character-matched control cannot isolate a cognitive anchoring mechanism.

The primary state-minus-feedback utility difference is
\$\PrimaryEstimate{} per failed task, with descriptive 95\% interval
[\$\PrimaryLow, \$\PrimaryHigh]. State and feedback selectors solve 50/94 and
49/94 tasks at mean API costs \$0.00011585 and \$0.00012066. The source-selected
static rule solves 60/94 for \$0.00055505. Its state-minus-static utility
difference is $-\$0.00062463$, with interval
$[-\$0.00142948,+\$0.00010830]$. Lower state-policy cost therefore does not
establish equal reliability or a utility improvement. The added escalation
diagnostic solves 58/94 for \$0.00026911, but is not substituted for the primary.

Local prediction adds no established artifact-feature benefit. State-minus-
feedback Brier difference is $-0.000358$, interval $[-0.007756,+0.006195]$;
the benefit-MSE difference is $+0.001937$, interval $[-0.002906,+0.007326]$.
All views' absolute losses are in Table~\ref{tab:mbpp-prediction}. The state
view's Brier score (0.471100) is much worse than the global prior's (0.279996).
The prespecified shift diagnostic finds 60/94 target prompts outside the source
character-length range and all 94 outside its line-count range. This diagnoses
a transfer limitation of these fitted structural baselines, without proving
that semantic artifact value is intrinsically unpredictable.

The two alternate draw rotations also leave the primary utility interval
spanning zero: estimates are \$0.00021320 for draws (1,2) and \$0.00011231 for
(2,0), compared with \$0.00011119 for (0,1). An exploratory noise-adjusted
heterogeneity diagnostic is positive for all four model/benefit combinations.
Both-minus-feedback latent-variance estimates are 0.1057 for 9B and 0.0777 for
27B, with descriptive intervals [0.0619,0.1474] and [0.0330,0.1279]. Under
conditionally independent draws, this suggests state variation beyond simple
three-draw noise; it does not show that the available predictors capture it.

\subsection{Semantic forecast control}
\begin{table*}[t]
\centering\small
\caption{Prespecified API forecast control, restricted to tasks with all
three forecast views and all thirty recovery observations. Success and cost
use two-draw replay at \$0.01, including the producer and forecaster fees.
Cost is USD $\times1{,}000$ per failed task.}
\label{tab:semantic}
\input{generated/semantic-results.tex}
\end{table*}
All 288 forecasts parse correctly and incur \$0.0285302, without additional
uncertain admissions. Prediction and policy comparisons use the same 94 complete
recovery states (Table~\ref{tab:semantic}); the two incomplete solver states
remain excluded even though all their forecasts exist.

The full-state API view has a worse Brier score than feedback-only:
0.329437 versus 0.291994, difference $+0.037443$ with descriptive 95\% interval
$[+0.013180,+0.062739]$. Its benefit-MSE difference is $+0.008188$, interval
$[-0.011756,+0.027600]$. At the \$0.01 ceiling, full-state replay solves 49/94
tasks versus feedback-only's 46/94, at mean costs \$0.00047711 and
\$0.00050439. Their utility difference is \$0.00034643, interval
$[-\$0.00030522,+\$0.00102463]$; this does not establish an artifact-information
gain. Against the source-selected static rule on those same tasks, API state
utility is lower by \$0.00109227, interval
$[-\$0.00209206,-\$0.00012983]$. These secondary, complete-case comparisons are
descriptive and unadjusted, not a replacement primary or a result about every
possible semantic predictor.

The mean full-state forecast fee is \$0.00011077, about 2.05 times the mean
single fresh 9B solver draw (\$0.00005403) in this bank. Its full fee is paid
before choosing a solver, even when the resulting forecast is unhelpful.

Primary collection costs \$0.83799455 in known charges and retains
\$0.0145152 in reservations. Including the \$0.5149634 prior campaign upper
amount and the semantic control, cumulative charges plus all reservations are
\$1.39600335, below the unchanged \$2 ceiling. The cost of collecting this bank
is separate from the per-task replay bills and does not claim deployment
outage costs are identified.

\section{Interpretation and scope}

Artifact reuse and artifact-aware selection are separate interventions.
A positive both-versus-fresh difference can arise from feedback without any
benefit from retaining code. A selector's advantage over an unconditional prior
can arise from model choice without demonstrating that it knows which work to
preserve. We therefore require the feedback-only and model-only comparisons
before attributing a result to the failed artifact.

A forecasting system must also pay for its decision. In the source study,
the mean state-forecast fee is about 1.7 times the price of a fresh 9B recovery
draw. Better success estimates would therefore need a decision benefit large
enough to cover that overhead. Conversely, a static rule selected for the actual
two-draw objective can match a learned selector even when a weaker static
baseline makes the latter look beneficial.

The complete-case context interaction and the uncertain primary decision gain
answer different questions. This particular 9B solver suffers more from adding
failed code than the 27B solver, while a frozen selector with cheap artifact
features does not demonstrate that it can identify the right action per task.
The source-trained global and static controls remain essential: the richer
feature models transfer poorly, and a lower monetary cost alone does not imply
an improved reliability--cost tradeoff.

The positive exploratory heterogeneity diagnostic leaves open a useful research
question, rather than closing it with a universal fresh-generation rule.
Context has different observed value across states, but learning that value
economically requires another task-disjoint test. Every MBPP+ result in this
report is now observed; using these tasks to design a new predictor would make
them development data. A later study should compare stronger semantic and
recovery-routing systems on a fresh task family, retain feedback-only and
strong static controls, and measure genuine multi-step child-artifact handoff.

\section{Limitations}

This study evaluates one visible checkpoint on small Python functions.
Two attempts reuse the original parent context; they do not implement or
validate a general multi-step agent controller. A different model from the
producer performs recovery, so this is not a controlled same-model self-repair
study. Provider versions and quantization may affect behavior, and model sizes
are not randomized interventions.

The public feedback is a particular regime of base inputs and acceptable outputs,
not complete execution traces, all actual wrong outputs, or a repository's test
suite. Public passes can be false positives, and full-suite success remains
oracle-relative. Reference compatibility and public-example audits do not
establish semantic correctness for every specification. Pretraining overlap or
semantic similarity between these public benchmarks cannot be ruled out.

Three draws per condition do not reveal a state's true individual context effect.
One failed artifact per task does not isolate within-task variation across
different failed programs. Structural ridge baselines and one zero-shot 9B
forecaster do not exhaust learnable semantic predictors. Failure to improve with
these baselines would not prove artifact value intrinsically unpredictable.
Task-level intervals rely on the declared sampling and dependence assumptions,
and small effects may remain unresolved. The two permanently missing recovery
cells prevent the frozen full-bank confirmatory analysis. Complete-case results
may be sensitive to censoring and to the disclosed continuation after a partial
audit; all-selected-state bounds preserve this uncertainty without asserting
missing at random.

The distractor is character-matched rather than token-matched. Its effect can
therefore combine content relevance, tokenization and packaging. Hard admission
uses deliberately conservative token ceilings; the resulting budget frontier
is not an optimal tokenizer-aware deployment frontier. API charges exclude
local compute and hardware opportunity cost. Unknown service admissions remain
reserved collection overhead, so complete deployment outage/retry cost is not
identified. Semantic forecast responses are collected after outcomes despite
their earlier-frozen, outcome-blind requests.

\section{Reproducibility}

The repository releases task/split hashes, frozen weights, public checkpoints,
matched requests and responses, predictions, reservation and settlement journals,
transport censoring, finite replay and plotting code. Offline reproduction
requires neither credentials, model calls, candidate execution nor GPUs.
Numerical refits tolerate floating-point platform roundoff; policy replay uses
the saved predictions. Transport continuation imports known observations once
and preserves every unknown admission rather than silently rerunning it.

The canonical bundle is
\path{experiments/mbpp-plus/context-value-transfer/results/replication-20261022}
in \url{https://github.com/EmaRimoldi/how-to-pick-a-model}. Its checksums cover
the raw journal, archived segment history, provenance, and saved analysis.
Source snapshots preserve the original freeze, the actual prefix collector,
both actual recovery collectors, the semantic collector and their shared
helpers. Each snapshot records whether its anchor is a directly recorded
runtime hash or a frozen Git commit; inert \texttt{.py.txt} copies are evidence,
not extra runnable entry points.

The separate paper builder, \path{docs/paper/recovery-value/build.py}, audits the
primary and semantic journals, requires the recomputed report to match the saved
analysis up to floating-point roundoff, regenerates tables and figures, and
compiles this manuscript in a fresh output directory. Its table generator
requires an explicit descriptive mode for incomplete banks and prints selected
coverage. The existing \texttt{main.tex}, \texttt{arxiv.tex} and theory anchor
remain separate research artifacts.

\newpage
\begin{thebibliography}{11}
\bibitem{coderescue} Q. He et al. \emph{CodeRescue: Budget-Calibrated Recovery
Routing for Coding Agents}. arXiv:2607.19338v1, 2026.
\url{https://arxiv.org/html/2607.19338v1}. The later withdrawal is administrative.
\bibitem{falsification} M. Iscan. \emph{Falsification, Not Exposure}.
arXiv:2606.31511, 2026. \url{https://arxiv.org/abs/2606.31511}.
\bibitem{darc} P. Wang et al. \emph{Diagnosis Before Recovery: Turning Agent
Failures into Selective Self-Correction}. 2026.
\url{https://arxiv.org/abs/2608.11772}.
\bibitem{ddbench} Y. Yan et al. \emph{Evaluating Agentic Code Repair
Capabilities in Distributed Systems}. 2026.
\url{https://arxiv.org/abs/2608.14863}.
\bibitem{selfrepair} T. X. Olausson et al. \emph{Is Self-Repair a Silver Bullet
for Code Generation?} ICLR, 2024. \url{https://arxiv.org/abs/2306.09896}.
\bibitem{selfdebug} X. Chen et al. \emph{Teaching Large Language Models to
Self-Debug}. \url{https://arxiv.org/abs/2304.05128}.
\bibitem{reforc} \emph{Re-FORC: Adaptive Reward Prediction for Efficient
Chain-of-Thought Reasoning}. \url{https://arxiv.org/html/2511.02130v2}.
\bibitem{aifunctions} Strands Labs. \emph{AI Functions economics}.
\url{https://github.com/strands-labs/ai-functions/blob/8f1485c2b685c83560bd91a0569cdf3af3176e75/docs/economics.md}.
\bibitem{evalplus} J. Liu, C. S. Xia, Y. Wang and L. Zhang.
\emph{Is Your Code Generated by ChatGPT Really Correct?}
NeurIPS, 2023. \url{https://arxiv.org/abs/2305.01210}.
\bibitem{mbpp} J. Austin et al. \emph{Program Synthesis with Large Language
Models}. 2021. \url{https://arxiv.org/abs/2108.07732}.
\bibitem{openrouter} OpenRouter. \emph{Zero-completion insurance and pricing}.
\url{https://openrouter.ai/blog/announcements/never-pay-for-empty-ai-responses-again/};
\url{https://openrouter.ai/pricing}.
\end{thebibliography}

\clearpage
\onecolumn
\appendix
\section{Full descriptive comparison tables}

\begin{table}[!ht]
\centering\small
\caption{All twelve MBPP+ context comparisons in percentage points.
The ordinary 95\% and nominal Bonferroni-adjusted 99.583\% intervals share
the same task clusters. The latter correct the declared family, while all
complete-case intervals remain descriptive under censored collection.}
\label{tab:all-context-effects}
\input{generated/context-effects.tex}
\end{table}

\begin{table}[!ht]
\centering\small
\caption{Paired state-minus-control utility on MBPP+ at the \$0.01 ceiling,
original draws 0 and 1. Units are USD $\times1{,}000$ per failed task.
Only the feedback comparison is primary; escalation is an added diagnostic.}
\label{tab:all-policy-comparisons}
\input{generated/policy-comparisons.tex}
\end{table}

\begin{table}[!ht]
\centering\small
\caption{All eight fixed actions on MBPP+ at the \$0.01 ceiling. Cost and
utility are USD $\times1{,}000$ per failed task, including producer cost.
These target outcomes do not select the source-frozen static baseline.}
\label{tab:all-fixed}
\input{generated/all-fixed-policies.tex}
\end{table}

\clearpage
\twocolumn
\section{Noisy benefits and finite decision horizons}

\subsection{Paired benefit prediction losses}

Let $D$ be an unbiased noisy estimate of a fixed state's true context benefit
$\delta$, and let $h_1,h_0$ be predictions fixed independently of the branch
outcomes conditional on that state. A standard squared-error identity gives
\begin{align*}
 &\mathbb E[(h_1-D)^2-(h_0-D)^2\mid s]\\
 &\qquad=(h_1-\delta)^2-(h_0-\delta)^2.
\end{align*}
The common sampling-noise contribution cancels in the paired loss difference.
This explains why we report paired prediction-risk differences while refusing
to treat a three-draw observed effect as a state's ground truth. It does not
remove sampling uncertainty, establish generalization beyond the task
population, or make absolute observed-benefit MSE a noise-free quantity.

For the exploratory heterogeneity diagnostic, write $D_{ir}$ for draw $r$'s
both-minus-feedback (or both-minus-fresh) difference in state $i$, with $R=3$.
Under conditional independence of repeated draws and independence of states,
the cross-product estimator
\begin{align*}
\widehat V={}&\frac1n\sum_i\frac{\sum_{r\ne t}D_{ir}D_{it}}{R(R-1)}\\
 &-\frac{\sum_{i\ne j}\overline D_i\overline D_j}{n(n-1)}
\end{align*}
targets between-state variance of the latent benefit. A finite-sample estimate
can be negative; we retain that value and its bootstrap interval. This is a
diagnostic application of cross-products, not a new estimator or evidence that
the remaining heterogeneity is predictable from the available features.

\subsection{Why a one-draw static rule can be weak at two draws}

Consider an illustrative fixed action with full-suite pass probability $p$,
public false-positive probability $q$, and public failure probability
$r=1-p-q$. Suppose two independent calls from the same original checkpoint are
always admissible, have mean cost $c$, and execution stops on a public pass.
Then
\[
 P(\text{full success by two draws})=p(1+r),
\]
and expected inference cost is $c(1+r)$. For solution value $v$ and common
prefix cost $k_0$, expected net utility is
\[
 (vp-c)(1+r)-k_0.
\]
Ranking actions by the one-draw score $vp-c$ need not rank them by this two-draw
objective: public false positives and the chance of admitting a second attempt
matter. Real hard-budget admission and variable costs add further constraints;
we therefore use recorded source replay rather than this simplifying formula
to select the static comparator. The calculation supplies no optimality
guarantee for the learned greedy selector or a multi-step agent.

\subsection{A limited decision-risk bound}

For a binary context choice with known incremental cost $\Delta k$, let
$g=v\delta-\Delta k$ and $\widehat g=v\widehat\delta-\Delta k$.
The oracle chooses the added context exactly when $g>0$, and a predictor chooses
it when $\widehat g>0$. Pointwise decision regret satisfies
\[
 |g|\,\mathbf1\{\mathbf1[g>0]\ne\mathbf1[\widehat g>0]\}
 \le v|\delta-\widehat\delta|.
\]
On a wrong-sign choice, the estimation error must span zero and hence exceed
$|g|$; on a correct choice regret is zero. Cauchy--Schwarz then gives the
illustrative bound $\mathbb E[\mathrm{regret}]\le
v\sqrt{\mathbb E[(\delta-\widehat\delta)^2]}$.
Cost-estimation error contributes its own term when costs are unknown. This
binary, known-cost statement does not imply a measured gain for our eight-action,
budget-constrained policies, whose costs and stopping behavior are empirical.


\end{document}
